\title{DeepSeekMath: Pushing the Limits of Mathematical Reasoning in Open Language Models}

\begin{abstract}
Mathematical reasoning remains a formidable obstacle for language models, owing to its inherently intricate and structured characteristics. This study presents DeepSeekMath~7B, a model developed by extending the pre-training of DeepSeek-Coder-Base-v1.5 7B using 120B mathematics-focused tokens gathered from Common Crawl, supplemented with data in both natural language and code. DeepSeekMath~7B achieves a remarkable 51.7\% score on the advanced MATH benchmark without dependence on external toolkits or ensemble-based voting methods, narrowing the gap with leading systems such as Gemini-Ultra and GPT-4. When applying self-consistency with 64 samples, DeepSeekMath~7B further raises its MATH benchmark score to 60.9\%. The advancements in DeepSeekMath~'s mathematical reasoning stem from two principal strategies: firstly, the effective utilization of extensive, publicly accessible web resources via a carefully crafted data selection process; and secondly, the development of Group Relative Policy Optimization (GRPO), an adaptation of Proximal Policy Optimization (PPO), which simultaneously boosts mathematical reasoning performance and improves PPO’s memory efficiency.
\end{abstract}

\section{Introduction}

Large language models (LLM) have transformed the landscape of mathematical reasoning within artificial intelligence, leading to notable progress on quantitative reasoning benchmarks \citep{MATH} and geometry reasoning benchmarks \citep{trinh2024solving}. Additionally, these models have become valuable tools for aiding humans in tackling intricate mathematical challenges \citep{tao}. Despite these achievements, leading systems like GPT-4 \citep{gpt4} and Gemini-Ultra \citep{gemini} remain closed to the public, while open-source alternatives available today substantially underperform in comparison.

This paper presents DeepSeekMath, a specialized language model designed for mathematics, which demonstrates superior mathematical reasoning compared to existing open-source models and delivers performance on par with GPT-4 across academic benchmarks. The foundation of this work is the DeepSeekMath Corpus, an extensive, high-quality pre-training dataset consisting of 120B math-related tokens. The corpus is assembled from the Common Crawl (CC) using a fastText-based classification method \citep{joulin2016fasttext}. Initially, the classifier is trained with positive samples sourced from OpenWebMath \citep{openwebmath}, supplemented by a varied range of negative examples drawn from unrelated web content. Following this, the trained classifier identifies further math-relevant data within CC, and these additional candidates are meticulously curated through manual review. The classifier is then retrained using this augmented dataset to boost its accuracy. Evaluation outcomes underscore the robustness of the resulting corpus, as DeepSeekMath-Base 7B achieves a 64.2\% score on GSM8K \citep{gsm8k} and 36.2\% on the high-difficulty MATH benchmark \citep{MATH}, thus exceeding the performance of Minerva 540B \citep{minerva}. Notably, because the DeepSeekMath Corpus encompasses multiple languages, we observe marked gains in Chinese mathematics evaluations \citep{wei2023cmath,agieval}. Our contributions to mathematical data curation offer a new foundation for ongoing developments in this area, suggesting considerable opportunity for further enhancements.

DeepSeekMath-Base is constructed upon DeepSeek-Coder-Base-v1.5 7B \citep{deepseek-coder}, as our observations suggest that initializing from a code-oriented pretrained model yields superior outcomes relative to those from a standard LLM. Additionally, we find that mathematical training not only strengthens mathematical proficiency, but also results in enhanced performance on benchmarks such as MMLU \citep{mmlu} and BBH \citep{bbh}, highlighting the benefit to broader reasoning skills.

Following the pre-training phase, we perform mathematical instruction tuning on DeepSeekMath-Base, leveraging datasets for chain-of-thought \citep{cot}, program-of-thought \citep{pot,pal}, and tool-integrated reasoning \citep{tora}. This instruction-tuned variant, DeepSeekMath-Instruct 7B, outperforms all other 7B models and demonstrates results comparable to 70B-scale open-source instruction-tuned models.

In addition, we present Group Relative Policy Optimization (GRPO), an alternative reinforcement learning (RL) approach derived from Proximal Policy Optimization (PPO) \citep{schulman2017proximal}. Unlike standard PPO, GRPO dispenses with the need for a critic model, estimating the baseline directly from group scores and thus substantially decreasing the computational demands of training. Using only a selected subset of English instruction tuning data, GRPO achieves notable performance gains over the already competitive DeepSeekMath-Instruct, both on in-domain benchmarks (e.g., GSM8K: 82.9\% $\rightarrow$ 88.2\%, MATH: 46.8\% $\rightarrow$ 51.7\%) and in out-of-domain settings (e.g., CMATH: 84.6\% $\rightarrow$ 88.8\%) during RL optimization. Additionally, we propose a comprehensive framework to interpret and relate multiple approaches, including Rejection Sampling Fine-Tuning (RFT) \citep{yuan2023scaling}, Direct Preference Optimization (DPO) \citep{dpo}, PPO, and GRPO. This unifying perspective reveals that each method operates as a variant or streamlined form of reinforcement learning. We also conduct a thorough set of empirical investigations—covering dimensions such as online versus offline training, outcome versus process-based supervision, and comparisons of single-turn with iterative RL—to elucidate the foundational characteristics of this unified scheme. Finally, we articulate how our RL methodology enhances the capabilities of instruction-tuned models, and discuss promising research avenues for designing even more efficient RL frameworks within this overarching paradigm.

\subsection{Contributions}
Our contribution includes scalable math pre-training, along with the exploration and analysis of reinforcement learning.

\textbf{Math Pre-Training at Scale}

\begin{itemize}[topsep=0pt]
    \item This study demonstrates that Common Crawl, an openly available resource, holds significant potential for mathematics-related tasks. Through the use of a carefully developed data selection pipeline, we are able to compile the DeepSeekMath Corpus—a comprehensive dataset comprising 120B tokens sourced from math-rich web pages. This corpus notably surpasses the scale of math web data used by Minerva \citep{minerva} by nearly sevenfold and is approximately nine times larger than the newly released OpenWebMath \citep{openwebmath}.

\item DeepSeekMath-Base 7B, our foundational pre-trained model, attains results on par with Minerva 540B \citep{minerva}, suggesting that parameter count alone does not determine mathematical reasoning proficiency. High performance can also be realized by a more compact model that is trained on superior quality data.

\item We present results from our experiments on math training. Pre-training models with code enhances their performance on mathematical problem solving, irrespective of whether external tools are employed. This finding contributes to addressing the persistent inquiry: \textit{does code training enhance reasoning skills?} Our evidence suggests it does, particularly in the context of mathematical reasoning.

    \item Although training on arXiv papers is common, especially in many math-related papers, it brings no notable improvements on all mathematical benchmarks adopted in this paper.
\end{itemize}

\textbf{Exploration and Analysis of Reinforcement Learning}

\begin{itemize}[topsep=0pt]
    \item We present Group Relative Policy Optimization (GRPO), a reinforcement learning algorithm that is both resource-efficient and performant. Unlike Proximal Policy Optimization (PPO), GRPO omits the use of a critic, utilizing group-based score estimation as the baseline, which leads to a notable reduction in the computational resources needed for training.
    \item We empirically show that DeepSeekMath-Instruct, when instruction-tuned and further optimized using GRPO solely with instruction-tuning data, achieves substantial performance gains. Moreover, our approach facilitates improved generalization to out-of-domain tasks within the reinforcement learning phase.
    \item We offer a unified analytical framework encompassing various approaches such as RFT, DPO, PPO, and GRPO. Through systematic experimentation—covering topics like online versus offline learning, process supervision compared to outcome supervision, and the differences between single-step and multi-step reinforcement learning—we thoroughly dissect the core constituents of this framework.
    \item Leveraging insights from our unified analytical paradigm, we investigate the underlying factors that drive the success of reinforcement learning. We also articulate several promising avenues for advancing the efficacy of reinforcement learning in large language models (LLMs).
\end{itemize}

\subsection{Summary of Evaluations and Metrics}

\begin{itemize}[topsep=0pt]
    \item \textbf{English and Chinese Mathematical Reasoning}: Our models undergo thorough evaluation on both English and Chinese datasets that span a spectrum of mathematical challenges from elementary to university levels. The English benchmarks utilized are GSM8K \citep{gsm8k}, MATH \citep{MATH}, SAT \citep{llemma}, OCW Courses \citep{minerva}, and MMLU-STEM \citep{mmlu}. For the Chinese language, our evaluation incorporates MGSM-zh \citep{mgsm}, CMATH \citep{wei2023cmath}, Gaokao-MathCloze \citep{agieval}, and Gaokao-MathQA \citep{agieval}. We measure the capability of the models both in producing complete textual solutions independently of external tools and in resolving problems by leveraging Python programming.

On English evaluation datasets, DeepSeekMath-Base demonstrates performance on par with the proprietary Minerva 540B model \citep{minerva}, and consistently outperforms all open-source base models—including Mistral 7B \citep{mistral} and Llemma-34B \citep{llemma}—regardless of their use of mathematical pre-training, frequently by a notable margin. Particularly, DeepSeekMath-Base excels on Chinese assessment benchmarks. This is likely attributable to our divergence from previous methodologies \citep{minerva,llemma}, as we include not only English-language but also high-quality mathematical data from other languages during pre-training. Incorporating mathematical instruction fine-tuning as well as reinforcement learning leads to the enhanced models DeepSeekMath-Instruct and DeepSeekMath-RL, which, for the first time in the open-source sphere, achieve over 50\% accuracy on the challenging, competition-grade MATH dataset.

\item \textbf{Formal Mathematics}: We assess DeepSeekMath-Base using the informal-to-formal theorem proving challenge introduced in \citep{dsp_proof}, employing miniF2F \citep{minif2f} and designating Isabelle \citep{isabelle} as the proof assistant for evaluation. DeepSeekMath-Base exhibits notable proficiency in few-shot autoformalization tasks.
\item \textbf{Natural Language Understanding, Reasoning, and Code}: For a holistic assessment of the models’ general comprehension, logical reasoning, and programming abilities, DeepSeekMath-Base is evaluated using the Massive Multitask Language Understanding (MMLU) benchmark \citep{mmlu}, which features 57 multiple-choice tasks from a wide array of disciplines, as well as BIG-Bench Hard (BBH) \citep{bbh}, comprising 23 complex tasks that often involve multi-step reasoning, in addition to HumanEval \citep{codex} and MBPP \citep{mbpp}, both standard measures for code language model evaluation. Pre-training on mathematical data enhances both language understanding and reasoning capabilities.

\section{Math Pre-Training}

\subsection{Data Collection and Decontamination} This section details the methodology used to assemble the DeepSeekMath Corpus utilizing data from Common Crawl. Figure \ref{fig:data_collect} illustrates an iterative workflow employed to efficiently compile an extensive mathematics-focused corpus from Common Crawl, beginning with a seed dataset—typically a compact, high-caliber set of math-related data. It should be emphasized that this pipeline is adaptable and can be utilized for corpus construction in various fields, including coding.

To begin, we select OpenWebMath \citep{openwebmath}, a curated repository of high-quality mathematical web materials, as the foundation for our initial corpus. Leveraging this corpus, we utilize a fastText model \citep{joulin2016fasttext} to identify additional web pages that are similar in nature to those in OpenWebMath. For the training set, we sample 500,000 entries at random from OpenWebMath to serve as positive examples, alongside 500,000 negative examples randomly drawn from Common Crawl. The fastText model is trained using an open-source toolkit\footnote{\url{https://fasttext.cc}}, with hyperparameters set as follows: a vector size of 256, a learning rate of 0.1, a maximum word n-gram length of 3, a minimum frequency threshold of 3, and 3 training epochs. To scale down the Common Crawl dataset, we apply URL-based deduplication and near-deduplication approaches, yielding a total of 40B unique HTML documents. Using the trained fastText classifier, we then retrieve mathematical web content from this deduplicated dataset. In order to eliminate lower-quality examples, we rank all retrieved pages based on their fastText-assigned scores, retaining only those that appear highest on this scale. We gauge how much data is preserved by conducting pre-training trials with the uppermost 40B, 80B, 120B, and 160B tokens. For this initial phase, the selection is restricted to the top 40B tokens.

Following the initial phase of data collection, a significant portion of mathematical web pages remains uncaptured, largely because the existing fastText model has been trained on positive examples that are not adequately varied. To address this, we seek out new mathematical web sources to expand and diversify the seed corpus, thus refining the fastText model. Our methodology begins with partitioning the entirety of Common Crawl into distinct domains, where each domain comprises web pages sharing the same base URL. We then determine the fraction of each domain's web pages that were obtained during the first iteration. If a domain has more than 10\% of its pages collected, it is flagged as math-related (for example, \url{mathoverflow.net}).

Next, we perform manual annotation of those URLs within these selected domains that pertain to mathematical content (such as \url{mathoverflow.net/questions}). Previously uncollected web pages associated with these annotated URLs are then incorporated into the seed corpus. This process allows us to curate a broader set of positive examples, resulting in a more robust fastText model that can identify even more mathematical web pages in the subsequent round. Upon completing four rounds of this procedure, our collection comprises 35.5M mathematical web pages, summing to 120B tokens. Notably, in the fourth round, we observe that nearly 98\% of web pages were already retrieved during the third round, prompting us to discontinue further data acquisition.

In order to prevent benchmark contamination, we adopt the approach described in \cite{deepseek-coder}, filtering out web pages that feature questions or answers from English mathematical benchmarks such as GSM8K \citep{gsm8k}, MATH \citep{MATH}, as well as Chinese datasets like CMATH \citep{wei2023cmath} and AGIEval \citep{agieval}. Our filtering process entails removing any text fragments from the math training data that include a 10-gram sequence that perfectly matches any substring found in these evaluation benchmarks. For benchmark samples whose total length is under 10 grams but at least 3 grams, we apply strict exact matching to eliminate any potentially contaminated content from the corpus.

\subsection{Validating the Quality of the DeepSeekMath~Corpus}

We run pre-training experiments to investigate how the DeepSeekMath~Corpus is compared with the recently released math-training corpora:

\begin{itemize}[topsep=0pt]
    \item \textbf{MathPile} \citep{mathpile}: a dataset compiling 8.9B tokens from a variety of sources, including textbooks, Wikipedia, ProofWiki, CommonCrawl, StackExchange, and arXiv. Notably, arXiv accounts for more than 85\% of the material;
    \item \textbf{OpenWebMath} \citep{openwebmath}: a collection comprised of 13.6B tokens drawn from CommonCrawl, with data specifically curated to emphasize mathematical content;
    \item \textbf{Proof-Pile-2} \citep{llemma}: an aggregate mathematical dataset formed by combining OpenWebMath, AlgebraicStack (containing 10.3B tokens of mathematical code), and arXiv papers (amounting to 28.0B tokens). For experiments involving Proof-Pile-2, we adhere to the protocol in \cite{llemma}, employing a source ratio for arXiv:Web:Code of 2:4:1.
\end{itemize}

\subsubsection{Training Setting}
\label{sec:quality-policy}
We fine-tune a 1.3B-parameter general-purpose language model, structured identically to the DeepSeek LLMs \citep{deepseek-llm} and referenced as DeepSeek-LLM 1.3B, for mathematical tasks. Each mathematical dataset serves as the basis for a distinct model, with training conducted over 150B tokens per corpus. The training procedures employ the resource-efficient HAI-LLM framework \citep{haillm}. Consistent with DeepSeek LLMs training methodology, we utilize the AdamW optimizer \citep{adamW} configured with $\beta_1=0.9$, $\beta_2=0.95$, and $\mathrm{weight\_decay}=0.1$. The learning rate is managed via a multi-phase scheduling strategy: after a warmup period of 2,000 steps, the learning rate peaks, then drops to 31.6\% of its maximum by 80\% completion of training, and further decays to 10.0\% of its peak by 90\% completion. The upper bound for the learning rate is set at 5.3e-4, utilizing batches containing 4 million tokens, with a maximum sequence context of 4K tokens.

\subsubsection{Evaluation Results}

\textbf{The DeepSeekMath~Corpus is of high quality, covers multilingual mathematical content, and is the largest in size.}

\begin{itemize}[topsep=0pt]
    \item \textbf{High-quality}:
    Downstream results on eight mathematical benchmarks are assessed using few-shot chain-of-thought prompting following \cite{cot}. According to Table \ref{tab:corpora_comparison}, models trained on the DeepSeekMath~Corpus significantly outperform those using alternative corpora. Further, Figure \ref{fig:corpora_comparisons} illustrates that the DeepSeekMath Corpus consistently yields superior results relative to Proof-Pile-2 at the 50B token mark (equivalent to one full epoch over Proof-Pile-2), supporting the notion that DeepSeekMath Corpus possesses a higher average data quality.
    \item \textbf{Multilingual}:
    The DeepSeekMath~Corpus contains a broad linguistic diversity, with English and Chinese as the most prominent languages represented. Table \ref{tab:corpora_comparison} demonstrates that utilizing DeepSeekMath~Corpus for model training leads to enhanced mathematical reasoning abilities in both English and Chinese. By contrast, prior mathematical corpora, largely centered on English content, fail to deliver substantial gains in Chinese and may even impede performance on Chinese mathematical reasoning benchmarks.
    \item \textbf{Large-scale}:
    Compared to other available mathematical corpora, the DeepSeekMath~Corpus offers substantially greater volume. As indicated in Figure \ref{fig:corpora_comparisons}, DeepSeek-LLM 1.3B, trained on DeepSeekMath~Corpus, experiences a notably accelerated and sustained improvement trajectory. Conversely, baseline corpora, which are considerably smaller and have been cycled through multiple times during model training, result in rapid saturation of performance improvements.
\end{itemize}

\subsection{Training and Evaluating DeepSeekMath-Base 7B}

This section presents DeepSeekMath-Base 7B, a foundational model demonstrating robust capabilities in mathematical reasoning. The model is built upon DeepSeek-Coder-Base-v1.5 7B \citep{deepseek-coder} and has undergone training with a total of 500B tokens. The data utilized is composed of 56\% from the DeepSeekMath Corpus, 4\% from AlgebraicStack, 10\% sourced from arXiv, 20\% from Github code repositories, and the last 10\% consisting of natural language text drawn from the Common Crawl dataset in both English and Chinese. Training largely follows the approach outlined in Section \ref{sec:quality-policy}, with the exceptions of setting the peak learning rate to 4.2e-4 and employing a batch size of 10M tokens.

Our study offers a thorough evaluation of the mathematical proficiency of DeepSeekMath-Base 7B. Specifically, we examine the model’s capacity to independently generate complete mathematical solutions without assistance, address mathematical questions utilizing external tools, and perform formal theorem proving tasks. In addition to mathematical assessment, we further characterize the general competencies of the base model, analyzing its capabilities in natural language understanding, logical reasoning, and programming.

\paragraph{Mathematical Problem Solving with Step-by-Step Reasoning}

We assess the capabilities of DeepSeekMath-Base in addressing mathematical questions using few-shot chain-of-thought prompting \citep{cot} over eight distinct benchmarks available in both English and Chinese. These benchmarks span a range of areas, including quantitative reasoning—such as GSM8K \citep{gsm8k}, MATH \citep{MATH}, and CMATH \citep{wei2023cmath}—as well as multiple-choice tasks, represented by MMLU-STEM \citep{mmlu} and Gaokao-MathQA \citep{agieval}. Collectively, they encompass mathematical topics from basic to collegiate difficulty levels.

As illustrated in Table \ref{tab:base_math_cot}, DeepSeekMath-Base 7B demonstrates superior results over all eight evaluated benchmarks, outperforming other open-source base models such as the extensively adopted Mistral 7B \citep{mistral} and the more recent Llemma 34B \citep{llemma}, the latter of which has been trained on Proof-Pile-2 for enhanced mathematical capability \citep{llemma}. In particular, DeepSeekMath-Base achieves an absolute performance improvement exceeding 10\% compared to available open-source base models on the challenging, competition-grade MATH dataset. Furthermore, it even surpasses Minerva 540B \citep{minerva}—a closed-source base model with 77 times more parameters, built upon PaLM \citep{palm} and subjected to additional training using mathematical corpora.

\paragraph{Mathematical Problem Solving with Tool Use} Our assessment of program-based mathematical reasoning is conducted on the GSM8K and MATH benchmarks, employing few-shot program-of-thought prompting as described by \citet{pot,pal}. In this setting, the models are instructed to address each question by generating a Python script, leveraging external libraries such as \textit{math} and \textit{sympy} for handling complex calculations. The output produced by executing the script is then taken as the final answer. According to the results presented in Table \ref{tab:base_math_pot_proof}, DeepSeekMath-Base 7B achieves superior performance compared to the previous leading model, Llemma 34B.

\paragraph{Formal Mathematics}

The automation of formal proofs plays a significant role in improving both the precision and dependability of mathematical reasoning, as well as boosting productivity—an area that has attracted considerable interest in recent years. In this work, we assess the capabilities of DeepSeekMath-Base 7B for the task of informal-to-formal proving, as outlined by \citep{dsp_proof}. This task involves producing a formal proof from an informal description, its formal statement, and an informal argument. Our evaluation makes use of miniF2F \citep{minif2f}, a standard test set featuring Olympiad-level mathematical problems, for which we prompt the model to produce formal proofs in Isabelle using few-shot approaches. Consistent with the methodology described in \cite{dsp_proof}, the models are employed to draft proof sketches, while the standard automated prover Sledgehammer \citep{sledgehammer} is used to complete any omitted details. According to Table \ref{tab:base_math_pot_proof}, DeepSeekMath-Base 7B achieves notable results in the task of automating formal proofs.

\paragraph{Natural Language Understanding, Reasoning, and Code}

We assess the abilities of the model in natural language comprehension using the MMLU benchmark \citep{mmlu}, in reasoning with BBH \citep{bbh}, and in programming skills on the HumanEval \citep{codex} and MBPP \citep{mbpp} datasets. According to Table \ref{tab:understanding_reasoning_code}, DeepSeekMath-Base 7B achieves notable improvements on both the MMLU and BBH tasks compared to its earlier version, DeepSeek-Coder-Base-v1.5 \citep{deepseek-coder}, indicating that math-centric training yields substantial benefits for tasks involving understanding and reasoning. Furthermore, by retaining code-related data during continued training, DeepSeekMath-Base 7B preserves the coding task performance of DeepSeek-Coder-Base-v1.5 across both code evaluation benchmarks. In summary, DeepSeekMath-Base 7B delivers a marked performance increase over the more general Mistral 7B model \citep{mistral} on all three evaluated benchmarks in reasoning and programming.

\section{Supervised Fine-Tuning}

\subsection{SFT Data Curation}

A mathematical instruction-tuning dataset is developed, encompassing both English and Chinese tasks drawn from a range of mathematical domains and representing diverse levels of difficulty. Each problem is accompanied by a corresponding solution, provided in multiple formats including chain-of-thought (CoT) \citep{cot}, program-of-thought (PoT) \citep{pot,pal}, and tool-integrated reasoning format \citep{tora}. Altogether, the dataset consists of 776K training instances.

\begin{itemize}[topsep=0pt]
    \item \textbf{English mathematical datasets}: We provide tool-integrated annotated solutions for problems drawn from GSM8K and MATH, and further include a curated portion of MathInstruct \citep{MathInstruct} as well as the Lila-OOD training data \citep{lila}, where tasks are addressed through either CoT or PoT approaches. This collection encompasses a broad spectrum of mathematical disciplines, including algebra, probability, number theory, calculus, and geometry.
    \item \textbf{Chinese mathematical datasets}: Our resources comprise Chinese K-12 math questions across 76 distinct sub-topics—such as linear equations—with each solution detailed both in CoT style and in a tool-integrated reasoning format.
\end{itemize}

\subsection{Training and Evaluating DeepSeekMath-Instruct 7B}

This section presents DeepSeekMath-Instruct 7B, an extension of DeepSeekMath-Base that has been further optimized through mathematical instruction tuning. During training, examples are combined in random order until they fill a maximum context window of 4K tokens. The training process consists of 500 steps, with each batch comprising 256 examples, and utilizes a fixed learning rate set to 5e-5.

To assess mathematical proficiency, we examine model performance under both tool-free and tool-assisted settings. Evaluations are conducted across four quantitative reasoning datasets, encompassing both English and Chinese. Our results are compared against contemporary state-of-the-art models:

\begin{itemize}[topsep=0pt]
    \item \textbf{Closed-source models} comprise the following: (1) the GPT series, notably GPT-4 \citep{gpt4} and the GPT-4 Code Interpreter~\footnote{\url{https://openai.com/blog/chatgpt-plugins\#\#code-interpreter}} as their most advanced representatives, (2) Gemini Ultra and Pro \citep{gemini}, (3) Inflection-2 \citep{inflection-2}, (4) Grok-1~\footnote{\url{https://x.ai/model-card}}; as well as recent offerings from Chinese enterprises such as (5) Baichuan-3~\footnote{\url{https://www.baichuan-ai.com}}, and (6) the most recent GLM-4~\footnote{\url{https://open.bigmodel.cn/dev/api\#glm-4}}

    from the GLM family \citep{glm}.

These models serve broad applications, with the majority having been subjected to extensive alignment protocols.
    \item \textbf{Open-source models} encompass: general-purpose models such as (1) DeepSeek-LLM-Chat 67B \citep{deepseek-llm}, (2) Qwen 72B \citep{qwen}, (3) SeaLLM-v2 7B \citep{seallm}, and (4) ChatGLM3 6B \citep{chatglm3}. Additionally, models tailored for improved mathematical capabilities include: (5) InternLM2-Math 20B~\footnote{\url{https://github.com/InternLM/InternLM-Math}}, based on InternLM2 with subsequent mathematics-focused pretraining and instruction tuning; (6) Math-Shepherd-Mistral 7B, which leverages PPO training \citep{schulman2017proximal} applied to Mistral 7B \citep{mistral} using a process-supervised reward approach; (7) the WizardMath series \citep{wizardmath}, aimed at elevating mathematical reasoning in Mistral 7B and Llama-2 70B \citep{llama2} through evolve-instruct (instruction tuning using AI-generated instructions) and PPO training with data predominantly from GSM8K and MATH; (8) MetaMath 70B \citep{metamath}, consisting of Llama-2 70B fine-tuned on an expanded version of GSM8K and MATH; (9) ToRA 34B \cite{tora}, where CodeLlama 34B is fine-tuned to integrate external tools for mathematical reasoning; and (10) MAmmoTH 70B \citep{MathInstruct}, which involves instruction-tuning Llama-2 70B on MathInstruct data.

As illustrated in Table \ref{tab:sft_rl_math}, when evaluated without access to external tools, DeepSeekMath-Instruct 7B exhibits robust capability in sequential reasoning. Remarkably, on the MATH dataset at the competition level, this model outperforms every available open-source system as well as most proprietary counterparts (including Inflection-2 and Gemini Pro), exceeding them by a margin of at least 9\% absolute. This performance advantage holds even against much larger models (such as Qwen 72B) and those specifically fine-tuned with mathematics-oriented reinforcement learning strategies (for example, WizardMath-v1.1 7B). Although DeepSeekMath-Instruct achieves results comparable to the Chinese proprietary models GLM-4 and Baichuan-3 on the MATH benchmark, it still falls short relative to GPT-4 and Gemini Ultra.

When assessed in a context that permits both natural language reasoning and the use of programs as problem-solving tools, DeepSeekMath-Instruct 7B attains nearly 60\% accuracy on the MATH benchmark, outperforming all previously released open-source systems. On additional evaluation datasets, our model matches the performance of DeepSeek-LLM-Chat 67B—the earlier leading model despite its size being an order of magnitude greater.  
\section{Reinforcement Learning}

\subsection{Group Relative Policy Optimization} Following the Supervised Fine-Tuning (SFT) phase, reinforcement learning (RL) has demonstrated its capability to enhance the mathematical reasoning skills of LLMs~\citep{wang2023math,wizardmath}. Here, we present Group Relative Policy Optimization (GRPO), a novel RL method developed to achieve both efficiency and effectiveness.

\subsubsection{From PPO to GRPO} Proximal Policy Optimization (PPO) \citep{schulman2017proximal} represents a prevalent actor-critic reinforcement learning technique and is commonly employed for RL-based fine-tuning of LLMs \citep{ouyang2022training}. The primary mechanism of PPO involves maximizing the following surrogate objective to update the policy of LLMs:

\begin{equation}
\footnotesize
    \mathcal{J}_{PPO}(\theta) = \mathbb{E}{[q \sim P(Q), o \sim \pi_{\theta_{old}}(O|q)]} \frac{1}{|o|} \sum_{t=1}^{|o|} \min \left[ \frac{\pi_\theta(o_{t} | q, o_{<t})}{\pi_{\theta_{old}}(o_{t} | q, o_{<t})} A_{t}, \text{clip} \left( \frac{\pi_\theta(o_{t} | q, o_{<t})}{\pi_{\theta_{old}}(o_{t} | q, o_{<t})}, 1 - \varepsilon, 1 + \varepsilon \right)  A_{t} \right] ,
\end{equation}

Here, $\pi_{\theta}$ denotes the current policy model, whereas $\pi_{\theta_{old}}$ stands for the prior policy. The variables $q$ and $o$ correspond to questions and responses sampled from the question corpus and the earlier policy $\pi_{\theta_{old}}$, respectively. The parameter $\varepsilon$ serves as a clipping-related hyperparameter introduced in PPO to promote training stability. The advantage term $A_t$ is obtained via Generalized Advantage Estimation (GAE) \citep{gae}, which leverages the reward sequence $\{r_{\ge t}\}$ together with a trainable value function $V_{\psi}$. Consequently, PPO involves concurrent training of a value function alongside the policy model. To address the risk of over-optimization relative to the reward model, a common strategy is to introduce a per-token KL-divergence penalty between a reference model and the generated outputs at every token position \citep{ouyang2022training}; specifically,

\begin{equation}
    r_{t} = r_\varphi(q, o_{\le t}) - \beta \log\frac{\pi_{\theta}(o_{t}|q, o_{<t})}{\pi_{ref}(o_{t}|q, o_{<t})},
\label{eq:PPO-reward}
\end{equation}

where $r_\varphi$  is the reward model, $\pi_{ref}$ is the reference model, which is usually the initial SFT model, and $\beta$  is the coefficient of the KL penalty.

Since the value function in PPO is usually implemented as a separate model with a size comparable to that of the policy network, this introduces considerable overhead in terms of both memory and computation. Furthermore, within reinforcement learning training procedures, the value function is employed as a baseline to compute the advantage and thereby reduce variance. In the particular case of LLMs, reward signals are generally provided only for the final token through the reward model, making it challenging to train a value function that remains reliable at every token position. To mitigate these challenges, we introduce Group Relative Policy Optimization (GRPO), depicted in Figure \ref{fig:grpo}, which eliminates the necessity of learning an explicit value function as in PPO. Instead, GRPO utilizes the mean reward from several sampled responses—each generated for the same prompt—as the baseline. Specifically, for each question $q$, we generate a set of responses $\{o_1, o_2, \cdots, o_G\}$ using the prior policy $\pi_{\theta_{old}}$ and train the policy network to maximize the following objective:

\begin{equation}
\footnotesize
\begin{split}
    \mathcal{J}_{GRPO}(\theta) &= \mathbb{E}{[q \sim P(Q), \{o_i\}_{i=1}^G \sim \pi_{\theta_{old}}(O|q)]}  \\
    & \frac{1}{G}\sum_{i=1}^G\frac{1}{|o_i|} \sum_{t=1}^{|o_i|} \left\{ \min \left[ \frac{\pi_\theta(o_{i,t} | q, o_{i,<t})}{\pi_{\theta_{old}}(o_{i,t} | q, o_{i,<t})} \hat{A}_{i,t}, \text{clip} \left( \frac{\pi_\theta(o_{i,t} | q, o_{i,<t})}{\pi_{\theta_{old}}(o_{i,t} | q, o_{i,<t})}, 1 - \varepsilon, 1 + \varepsilon \right)  \hat{A}_{i,t} \right] - \beta \mathbb{D}_{KL}\left[\pi_{\theta} || \pi_{ref}\right]\right\} ,
\end{split}
\label{eq:GRPO-obj}
\end{equation}

The expected objective $\mathcal{J}_{GRPO}(\theta)$ is defined by first sampling $q$ from $P(Q)$ and a batch $\{o_i\}_{i=1}^G$ from the old policy $\pi_{\theta_{old}}(O|q)$. For each sequence $o_i$ in the batch, we compute an average over all time steps $t=1,\ldots,|o_i|$ as follows. The core inside the summation applies the minimum between the scaled advantage—where the ratio between the current and previous policy probabilities $\frac{\pi_\theta(o_{i,t} | q, o_{i,<t})}{\pi_{\theta_{old}}(o_{i,t} | q, o_{i,<t})}$ multiplies the estimated advantage $\hat{A}_{i,t}$—and a version where the ratio is constrained via the $\text{clip}$ operator to the interval $[1 - \varepsilon, 1 + \varepsilon]$. From this minimum, a regularization penalty $\beta \mathbb{D}_{KL}\left[\pi_{\theta} || \pi_{ref}\right]$ is subtracted, which encourages the learned policy to remain close to the reference policy $\pi_{ref}$. The full objective is the mean of these terms over all samples in the batch and over all steps within each sequence.

In this context, $\varepsilon$ and $\beta$ serve as hyper-parameters, while $\hat{A}_{i,t}$ denotes the advantage, which is computed solely based on the relative rewards among the outputs within each individual group; a more precise explanation is provided in the subsequent subsections. The strategy adopted by GRPO for advantage calculation—anchored in relative group rewards—coheres with the intrinsic comparative structure of reward models, since such models are customarily trained on comparison datasets involving multiple responses to the same prompt. Furthermore, rather than incorporating the KL penalty directly into the reward, GRPO imposes regularization by appending the KL divergence between the learned policy and the reference policy directly to the loss function, thereby simplifying the evaluation of $\hat{A}_{i,t}$. Unlike the KL penalty applied in (\ref{eq:PPO-reward}), GRPO estimates the KL divergence using the following unbiased estimator \citep{kl_approx}:

\begin{equation}
\small
    \mathbb{D}_{KL}\left[\pi_{\theta} || \pi_{ref}\right] = \frac{\pi_{ref}(o_{i,t}|q,o_{i,<t})}{\pi_{\theta}(o_{i,t}|q,o_{i,<t})}- \log\frac{\pi_{ref}(o_{i,t}|q,o_{i,<t})}{\pi_{\theta}(o_{i,t}|q,o_{i,<t})} - 1,
\end{equation}

which is guaranteed to be positive.

\begin{algorithm}[t]
  \small
  \caption{Iterative Group Relative Policy Optimization}
  \textbf{Input} starting policy model $\pi_{\theta_{\text{init}}}$; set of reward models $r_\varphi$; collection of task prompts $\mathcal{D}$;
  set hyperparameters $\varepsilon$, $\beta$, $\mu$
  \begin{algorithmic}[1]
    \State Initialize policy model: $\pi_\theta \leftarrow \pi_{\theta_{\text{init}}}$
    \For{iteration = 1, \dots, I}
      \State Set reference model: $\pi_{ref} \leftarrow \pi_{\theta}$
      \For{step = 1, \dots, M}
        \State Draw minibatch $\mathcal{D}_b$ from $\mathcal{D}$
        \State Store current policy: $\pi_{\theta_{old}} \leftarrow \pi_{\theta}$
        \State For each query $q \in \mathcal{D}_b$, sample $G$ responses $\{o_i\}_{i=1}^G \sim \pi_{\theta_{old}} (\cdot \mid q)$
        \State For each output $o_i$, calculate reward $r_i$ using $r_\varphi$
        \State For every token $t$ in $o_i$, estimate $\hat{A}_{i,t}$ via group relative advantage computation
        \For{GRPO iteration = 1, \dots, $\mu$}
          \State Optimize policy model $\pi_\theta$ by maximizing the GRPO objective (refer to Equation \ref{eq:GC-GRPO})
        \EndFor
      \EndFor
      \State Perform continual training on $r_\varphi$ using replay to update the reward model
    \EndFor 
  \end{algorithmic}
  \textbf{Output} $\pi_\theta$
  \label{alg:iter-grpo}
\end{algorithm}

\subsubsection{Outcome Supervision RL with GRPO} In this approach, given a question $q$, a set of candidate responses $\{o_1, o_2, \cdots, o_G\}$ is generated by sampling from the pre-existing policy model $\pi_{\theta_{old}}$. Each generated output is evaluated using a reward model, resulting in a set of $G$ reward scores $\mathbf{r}=\{r_1, r_2, \cdots, r_G\}$. These raw rewards are then standardized by subtracting the mean reward within the group and dividing by the group's standard deviation. Outcome supervision assigns this standardized reward to the terminal point of each output $o_i$ and propagates the normalized value as the estimated advantage $\hat{A}_{i, t}$ for all tokens within the respective output. Specifically, $\hat{A}_{i, t}$ is set to $\widetilde{r}_i = \frac{r_i- {\rm mean}(\mathbf{r})}{{\rm std}(\mathbf{r})}$. The policy parameters are subsequently updated by maximizing the objective presented in equation (\ref{eq:GRPO-obj}).

\subsubsection{Process Supervision RL with GRPO}

Outcome supervision delivers a single reward signal only after the generation of the complete output, which can be inadequate or inefficient for training policies in the context of intricate mathematical problems. Building on the approach in \cite{wang2023math}, we investigate process supervision, wherein a reward is allocated at the conclusion of each intermediate reasoning step. Specifically, for a given question $q$ and a set of $G$ sampled completions $\{o_1, o_2, \cdots, o_G\}$, a process reward model is deployed to evaluate each step across all sampled outputs, resulting in a set of rewards: $\mathbf{R} = \{ \{r_1^{index(1)},\cdots,r_1^{index(K_1)}\}, \cdots,  \{r_G^{index(1)},\cdots,r_G^{index(K_G)}\} \}$. Here, $index(j)$ corresponds to the position of the end token for the $j$-th reasoning step, and $K_i$ indicates the total number of steps in output $i$. These raw rewards are further standardized using their mean and standard deviation, that is, $\widetilde{r}_i^{index(j)} = \frac{r_i^{index(j)} - {\rm mean(\mathbf{R})}}{{\rm std(\mathbf{R})}}$. 

Thereafter, the advantage associated with each token is calculated as the aggregate of all normalized rewards from subsequent steps onward, mathematically denoted as $\hat{A}_{i, t} = \sum_{index(j) \ge t} \widetilde{r}_i^{index(j)}$. The policy is subsequently updated by maximizing the objective introduced in equation (\ref{eq:GRPO-obj}).

\subsubsection{Iterative RL with GRPO} During reinforcement learning, the existing reward model can become inadequate for effectively guiding the updated policy model. To address this limitation, we investigate iterative RL with GRPO. As illustrated in Algorithm \ref{alg:iter-grpo}, the iterative GRPO framework involves creating new datasets for the reward model, which are derived from the most recent samples generated by the policy model. The previous reward model is further trained through a replay mechanism that integrates 10\% of data from earlier iterations. Subsequently, the updated policy model is designated as the new reference model, and training of the policy model proceeds utilizing the newly refreshed reward model.

\subsection{Training and Evaluating DeepSeekMath-RL}

Our reinforcement learning experiments utilize the DeepSeekMath-Instruct 7B model. For RL training, we select chain-of-thought questions from the GSM8K and MATH subsets in the SFT dataset, totaling approximately 144K questions. To assess how RL influences benchmarks without direct exposure during fine-tuning, we deliberately omit other SFT questions. The construction of our reward model’s training set follows the approach of \citep{wang2023math}. We initialize the reward model using DeepSeekMath-Base 7B, adopting a learning rate of 2e-5. In the GRPO framework, the learning rate for the policy model is set to 1e-6, and we use a KL penalty coefficient of 0.04. Each question is associated with $64$ generated completions. We maintain a maximum sequence length of 1024 and set the batch size for training at 1024. The policy model undergoes a single update after every round of exploration. For evaluation, DeepSeekMath-RL 7B is assessed on the same benchmarks as DeepSeekMath-Instruct 7B. In the context of DeepSeekMath-RL 7B, chain-of-thought GSM8K and MATH tasks are categorized as in-domain, while all additional benchmarks are treated as out-of-domain.

Table \ref{tab:sft_rl_math} presents results for both open-source and closed-source models evaluated with chain-of-thought and tool-augmented reasoning strategies on English and Chinese test sets. Key observations include: 1) Using chain-of-thought reasoning, DeepSeekMath-RL 7B achieves 88.2\% and 51.7\% accuracy on GSM8K and MATH, respectively—outperforming every open-source model between 7B and 70B parameters, along with most closed-source counterparts. 2) Notably, DeepSeekMath-RL 7B's training is limited solely to chain-of-thought formatted instruction tuning datasets from GSM8K and MATH, and is initialized from DeepSeekMath-Instruct 7B. In spite of its restricted training data, DeepSeekMath-RL 7B delivers superior results across all considered metrics compared to DeepSeekMath-Instruct 7B, underlining the impact of reinforcement learning.

\section{Discussion}
In this section, we will share our findings in pre-training and RL experiments. 

\subsection{Lessons Learnt in Pre-Training}

To begin, we discuss our observations during the pre-training phase. Except where indicated, the training configurations described in Section \ref{sec:quality-policy} remain in effect. Notably, all references to the DeepSeekMath Corpus within this section pertain to an 89B-token subset obtained during the second round of data collection.

\subsubsection{Code Training Benefits Mathematical Reasoning}

An often-cited but not conclusively established hypothesis posits that code training enhances reasoning capabilities. In this work, we seek to partially address this claim, focusing on mathematics: code training enhances the capacity of models to perform mathematical reasoning, regardless of whether external tools are utilized.

To study how code training affects mathematical reasoning, we experimented with the following two-stage training and one-stage training settings:

\textbf{Two-Stage Training}

\begin{itemize}[topsep=0pt]
    \item \textbf{Code Training for 400B Tokens $\rightarrow$ Math Training for 150B Tokens}: Our approach involves pretraining DeepSeek-LLM 1.3B on 400B code tokens, then continuing training for an additional 150B tokens specifically focused on mathematical content.
    \item \textbf{General Training for 400B Tokens $\rightarrow$ Math Training for 150B Tokens}: For comparison, we perform a parallel experiment where, instead of code tokens, we utilize 400B general tokens—randomly drawn from a diverse large-scale dataset curated by DeepSeek-AI—in the initial pretraining phase, followed by the same 150B tokens of mathematical training, to assess the comparative benefits that code-specific pretraining may offer for mathematical reasoning capabilities.
\end{itemize}

\textbf{One-Stage Training}

\begin{itemize}[topsep=0pt]
    \item \textbf{Math Pretraining on 150B Tokens}: DeepSeek-LLM 1.3B is pretrained using 150B tokens exclusively from mathematical data.
    \item \textbf{Joint Training with 400B Code Tokens and 150B Math Tokens}: It has been observed that math finetuning after code pretraining negatively impacts coding abilities. To address this, we explore whether blending code and math data for concurrent training in a single phase can simultaneously enhance mathematical proficiency and reduce catastrophic forgetting.
\end{itemize}

\paragraph{Results}
Table \ref{tab:code-math} and Table \ref{tab:code-math-general-eval} demonstrate the downstream performance under different training settings.

Program-aided mathematical reasoning sees notable enhancement through code training, regardless of whether a two-stage or one-stage training strategy is employed. As evidenced in Table \ref{tab:code-math}, code training in the two-stage setup alone results in marked improvement in tackling GSM8K and MATH tasks using Python, and further gains are observed when math training is introduced in the subsequent stage. Notably, in the one-stage configuration, integrating both code tokens and math tokens helps prevent the catastrophic forgetting observed in two-stage approaches, while also providing mutual reinforcement for coding abilities (Table \ref{tab:code-math-general-eval}) and program-aided mathematical problem-solving (Table \ref{tab:code-math}).

Training with code also enhances mathematical reasoning abilities in settings where tools are not utilized. Within a two-stage training paradigm, introducing code training in the preliminary phase yields moderate gains by itself. Furthermore, this approach facilitates more effective learning during later math training, ultimately achieving superior outcomes. Conversely, when code and math tokens are integrated into a single-stage training process, the ability to reason mathematically without tools is diminished. This may be attributed to DeepSeek-LLM 1.3B’s relatively small model size, which restricts its ability to jointly learn and integrate both code and mathematical content during a single training stage.

\subsubsection{ArXiv Papers Seem Ineffective in Improving Mathematical Reasoning}

ArXiv papers frequently serve as a foundational element in math pre-training datasets \citep{minerva,gpt-f,llemma,mathpile}. Despite their widespread use, there has been limited comprehensive study into how they affect models' capabilities for mathematical reasoning. Somewhat unexpectedly, our experimental results indicate that arXiv papers do not substantially enhance mathematical reasoning skills. To investigate this, we tested a range of model architectures and scales, such as DeepSeek-LLM 1.3B and DeepSeek-Coder-Base-v1.5 7B \citep{deepseek-coder}, applying multiple preprocessing methodologies to the arXiv dataset:

\begin{itemize}[topsep=0pt]
    \item \textbf{MathPile} \citep{mathpile}: a dataset comprising 8.9 billion tokens, curated using a series of cleaning and filtering heuristics, with more than 85\% of its content sourced from scientific papers on arXiv;
    \item \textbf{ArXiv-RedPajama} \citep{redpajama}: a collection consisting of all arXiv LaTeX documents, where extraneous elements such as preambles, comments, macros, and bibliographies have been stripped, amounting to 28.0 billion tokens in total.
\end{itemize}

In the course of our study, we individually pretrain DeepSeek-LLM 1.3B on 150B tokens and DeepSeek-Coder-Base-v1.5 7B on 40B tokens sourced from each respective arXiv corpus. Our observations indicate that exposure to arXiv papers alone does not substantially enhance mathematical reasoning abilities. When utilizing a corpus comprised solely of arXiv data, neither model achieves measurable progress and, in some cases, even experiences a decline in performance across a suite of mathematical evaluation tasks of varying complexity featured in this work. These tasks encompass quantitative reasoning datasets such as GSM8K and MATH (Table \ref{tab:arxiv-cot}), multiple-choice evaluations like MMLU-STEM (Table \ref{tab:arxiv-cot}), as well as assessments of formal mathematics including miniF2F (Table \ref{tab:arxiv_atp}).

However, this conclusion has its limitations and should be taken with a grain of salt. We have not yet studied:

\begin{itemize}[topsep=0pt]
    \item How arXiv tokens influence specialized mathematical tasks not explored here, including the transformation of formal theorem statements or proofs into their informal counterparts;
    \item The interaction between arXiv tokens and additional data modalities;
    \item If advantages derived from arXiv papers would become more apparent as the model size increases.
\end{itemize}

Thus, further exploration is required, which we leave for future studies.

\subsection{Insights of Reinforcement Learning}
\subsubsection{Towards to a Unified Paradigm} In the following, we introduce a general framework that facilitates the comparative analysis of several training approaches, namely SFT, RFT, DPO, PPO, and GRPO, while also empirically examining the constituent elements of this unified paradigm. Typically, the parameter gradient $\theta$ for these methods can be uniformly expressed as:

\begin{equation}
    \nabla_{\theta}\mathcal{J}_{\textcolor{red}{\mathcal{A}}}(\theta) = \mathbb{E}[\underbrace{(q,o) \sim \textcolor{red}{\mathcal{D}}}_{Data \ Source}]\left( \frac{1}{|o|} \sum_{t=1}^{|o|}  \underbrace{GC_{{\mathcal{A}}}(q, o, t, \textcolor{red}{\pi_{{rf}}})}_{Gradient \ Coefficient}  \nabla_{\theta}\log \pi_{\theta}(o_t | q, o_{<t})\right).
\label{eq:objective}
\end{equation}

Three primary elements can be identified: 1) \textit{Data Source $\mathcal{D}$}, responsible for specifying the dataset used in training; 2) \textit{Reward Function $\pi_{{rf}}$}, providing the reward signals essential for guiding the training process; and 3) \textit{Algorithm $\mathcal{A}$}, which integrates the training dataset and reward information to generate the gradient coefficient $GC$ that modulates the strength of penalization or reward assigned to each data point. In the following, we examine a number of widely-used approaches within this general framework:

\begin{itemize}
    \item  \textbf{Supervised Fine-tuning (SFT)}: In SFT, a pretrained model undergoes additional training using SFT datasets curated by humans.
    \item \textbf{Rejection Sampling Fine-tuning (RFT)}: RFT enhances the SFT-trained model by conducting further fine-tuning on its responses, which are sampled from the model itself in response to SFT prompts. The selected outputs pass a filtering process that selects only those with correct answers.
    \item \textbf{Direct Preference Optimization (DPO)}: DPO incrementally improves the SFT model by tuning it with supplementary responses generated by the SFT model, leveraging pairwise DPO loss during training.
    \item \textbf{Online Rejection Sampling Fine-tuning (Online RFT)}: Unlike standard RFT, Online RFT begins with the SFT model as the initial policy model, subsequently updating it via fine-tuning with outputs generated dynamically from the evolving policy model.
    \item \textbf{PPO/GRPO}: In PPO/GRPO, the policy model is seeded with the SFT model's parameters and is further optimized using outputs sampled continuously from the live policy model during reinforcement learning.
\end{itemize}

We summarize the components of these methods in Table \ref{tab:data_policy}.
Please refer to Appendix \ref{app:analysis-rl} for a more detailed derivation process.

\paragraph{Observation about Data Source} We categorize data sources into two main types: online sampling and offline sampling. Online sampling refers to collecting training data generated by the current, actively exploring policy model during training, whereas offline sampling uses training data obtained from the sampling results of the original SFT model. Approaches like RFT and DPO utilize offline data, while methods such as Online RFT and GRPO rely on data gathered through online sampling.

As illustrated in Figure \ref{fig:rl-analysis}, our results indicate that Online RFT achieves noticeably better performance than RFT on both benchmarks. Initially, the performance of Online RFT closely matches that of RFT during the early phases of training; however, as training progresses, Online RFT obtains a marked lead, highlighting the benefits of online training methods. This result is expected: early in training, both the actor and the SFT model are highly similar, and the sampled data only reflect subtle distinctions. Later on, as the actor's parameters diverge further from the SFT model, the data generated by the actor become more distinct, making real-time sampling methods considerably more effective.

\paragraph{Observation about Gradient Coefficient} The algorithm utilizes the reward signal to modulate the gradient coefficient, which in turn governs the updating of the model parameters. In our experimental setup, we categorize the reward function as either `Rule' or `Model'. The `Rule' component assesses response quality through answer correctness, whereas the `Model' approach involves training a reward model to assign scores to responses, where the reward model itself is supervised using data labeled according to rule-based judgments. As shown in Equations \ref{eq:GC-RFT} and \ref{eq:GC-GRPO}, an important distinction arises between GRPO and Online RFT: only GRPO modifies its gradient coefficient according to the score received from the reward model. This enables GRPO to assign varying degrees of positive or negative feedback to responses contingent on the size of their respective rewards. Conversely, Online RFT does not possess this capacity—it applies no penalties for erroneous responses and grants the same level of reinforcement to all correctly answered responses, regardless of score.

As illustrated in Figure \ref{fig:rl-analysis}, GRPO outperforms online RFT, underscoring the effectiveness of adjusting the coefficients for positive and negative gradients. Moreover, the results demonstrate that GRPO+PS achieves better outcomes than GRPO+OS, suggesting that employing step-specific, fine-grained gradient coefficients is advantageous. Additionally, we investigate iterative RL within our experimental setup, carrying out two iterations. Figure \ref{fig:iter-rl} reveals that iterative RL leads to marked performance gains, most notably during the initial iteration.

\subsubsection{Why RL Works?} In this study, we apply reinforcement learning to a selected portion of the instruction tuning dataset and observe marked improvements over the baseline instruction tuning model. To delve deeper into the reasons behind RL’s effectiveness, we assess the Pass@K and Maj@K accuracies for both the Instruct and RL models using two benchmark datasets. As illustrated in Figure \ref{fig:maj-pass}, RL substantially increases Maj@K accuracy while leaving Pass@K largely unchanged. This pattern suggests that RL primarily strengthens the reliability of the model’s output distribution, \textbf{implying that the gains stem from elevating correct answers to the TopK outputs, as opposed to advancing the model’s intrinsic abilities.} Along similar lines, \citep{wang2023making} point out a \textbf{misalignment problem} in reasoning tasks within the SFT framework, revealing that SFT model reasoning skills can benefit from tailored preference alignment approaches \citep{yuan2023rrhf,song2023preference,wang2023making}.

\subsubsection{How to Achieve More Effective RL?}
\label{sec:effec_rl}
Our findings show that RL can be highly effective for tasks involving mathematical reasoning. Additionally, we outline a cohesive framework to interpret various well-known training methodologies, framing them as either variants of or approximations to RL. As indicated in Equation \ref{eq:objective}, the central elements of this paradigm include the Data Source, Algorithm, and Reward Function. Below, we suggest possible avenues for advancement concerning each of these three aspects.

\paragraph{Data Source} The data source constitutes the foundational input for all training methodologies. Within RL settings, we define the data source as consisting of unlabeled questions paired with outputs that are generated through sampling from the policy model. For the purposes of this study, we restrict ourselves to using questions gathered from the instruction tuning phase and employ a straightforward nucleus sampling approach for output generation. We hypothesize that this restriction may contribute to the observed improvement being limited to Maj@K within our RL pipeline. Moving forward, we intend to evaluate our RL framework on out-of-distribution prompts, and also integrate \textbf{advanced sampling (decoding) strategies}, such as those utilizing tree-search algorithms \citep{yao2023tree}. Furthermore, the application of \textbf{efficient inference techniques} \citep{xia-etal-2023-speculative,leviathan2023fast,kwon2023efficient,xia2024unlocking}, which impact the exploration capabilities of policy models, is recognized as critically influential.

\paragraph{Algorithms} Algorithms utilize both the available data and the received reward signal to compute the gradient coefficient, thereby facilitating updates to the model parameters. In accordance with Equation \ref{eq:objective}, contemporary approaches almost universally place full \textbf{TRUST} in the information provided by the reward function to modify the conditional likelihood of particular tokens. Yet, this reliance is problematic in practice because the reward signal may be unreliable, particularly in highly intricate tasks. For instance, even datasets like PRM800K \citep{lightman2023let}, despite rigorous annotation by skilled professionals, still include roughly 20\% erroneous annotations\footnote{\url{https://github.com/openai/prm800k/issues/12\#issuecomment-1728491852}}. In response, we turn our focus toward reinforcement learning algorithms specifically designed to withstand noisy reward feedback. It is our view that the adoption of \textbf{WEAK-TO-STRONG} \citep{burns2023weak} alignment strategies could profoundly influence and transform existing learning algorithms.

\paragraph{Reward Function} The reward function serves as the primary source of training signals in the reinforcement learning framework. Typically, in RL, a neural reward model fulfills this role. We identify three pivotal research directions concerning reward models: 1) \textbf{Strategies to improve the generalization capability of the reward model.} It is imperative for the reward model to generalize adequately to previously unseen queries and complex decoding outputs; failure to do so might result in RL simply preserving the LLMs’ existing output distribution instead of driving meaningful performance improvements; 2) \textbf{Methods to incorporate reward model uncertainty.} Understanding and representing the uncertainty within reward predictions can potentially connect the limitations of current reward models with weak-to-strong learning techniques; 3) \textbf{Techniques for constructing high-quality process reward models efficiently} that deliver granular feedback throughout the reasoning process \citep{lightman2023let,wang2023math}.

\section{Conclusion, Limitation, and Future Work} In this work, we introduce DeepSeekMath, an open-source model that surpasses all other open models on the competition-level MATH benchmark and narrows the gap with proprietary solutions. DeepSeekMath~begins with DeepSeek-Coder-v1.5 7B as its foundation, followed by extended training spanning 500B tokens, including a considerable fraction—specifically, 120B math-related tokens—extracted from Common Crawl. Through a comprehensive ablation analysis, we demonstrate that web-based sources provide substantial high-quality mathematical content, whereas arXiv's contribution falls short of initial expectations. We propose Group Relative Policy Optimization (GRPO), an adaptation of Proximal Policy Optimization (PPO), which enhances mathematical reasoning while reducing memory requirements. Experimental outcomes indicate that GRPO remains beneficial even when DeepSeekMath-Instruct 7B achieves strong benchmark results. Furthermore, we outline a unified framework for interpreting a variety of approaches and discuss several promising avenues for advancing reinforcement learning techniques.

While DeepSeekMath~demonstrates strong results on benchmarks assessing quantitative reasoning, its performance in geometry and theorem-proving tasks lags behind that of proprietary models. For example, during preliminary evaluations, the model was unable to solve questions involving triangles and ellipses, possibly reflecting a bias introduced during the data selection phase of pre-training and fine-tuning. Furthermore, due to its current scale, DeepSeekMath~does not match GPT-4’s proficiency in few-shot scenarios; GPT-4 benefits from few-shot prompts, whereas DeepSeekMath~shows little difference between its zero-shot and few-shot outcomes. Moving forward, we intend to enhance our data selection processes to develop a more robust and diverse pre-training corpus. Additionally, we plan to investigate avenues outlined in Section \ref{sec:effec_rl} to advance reinforcement learning techniques for large language models.

\end{document}